\section*{ID10036: Relation: -Plane) Like}
\paragraph{Metadata.}
PiID: 7445455129656 | RTEID: RTE:P38972.0168:T1.3820 | UUID: 66a17a59-6f13-5229-a8b4-07433260582b

\paragraph{Canonical Statement.}
Mathematically, we can idealize the local behavior near TZP as similar to a saddle point or a branch point. If we had to parameterize it, we might use something like lemniscatic coordinates or a figure-eight shaped potential. The lemniscate of Bernoulli (one form of a figure-eight curve) has an equation in the plane (say \$xy\$-plane) like \$(x\textasciicircum{}2 + y\textasciicircum{}2)\textasciicircum{}2 = a\textasciicircum{}2(x\textasciicircum{}2 - y\textasciicircum{}2)\$ 15 . It indeed goes through the origin \$(0,0)\$ as a crossing. If you revolve a lemniscate of Bernoulli around one of its symmetry axes, you get a sort of figure-eight shaped surface or a toroidal figure with a pinched middle. It’s plausible that the field lines in 3D near a dipole could be described by such a surface – for instance, if one considered lines of constant field magnitude or potential, you might get a lemniscate-like shape dividing regions of different

\paragraph{Canonical Equation.}
\[
-plane) like
\]

\paragraph{Dictionary.}
Mathematically, we can idealize the local behavior near TZP as similar to a saddle point or a branch point.

\paragraph{Encyclopedia.}
Relation: -Plane) Like is a symbolic expression, written -plane) like. It introduces a symbol or relation used elsewhere in the framework. It is built from p, l, a, n, k. Within the Trawin Zero Point registry it serves as one element of the framework's mathematical narrative.
